% ============================================================================
% Chapter 19 --- Wave 1: Identity and federation
% ----------------------------------------------------------------------------
% Purpose : align the identity layer with the contract of Ch. 6. First wave
%           with substantial trajectory differentiation. Sets the example
%           for Ch. 20-27.
% Page target: ~6 pages.
% Dependencies: Wave 0 (Ch. 18); reference layer Ch. 6.
% Mirrors the wave chapter template established in Ch. 18 (D18.2).
% ----------------------------------------------------------------------------
% Decisions registered in _coherence-EN.md:
%   D19.1 coexistence pattern: dual IdP for 12-24 months (typical)
%   D19.2 brownfield single-suite-centric trajectory: keep Entra, add OPA in front
%   D19.3 brownfield-hybrid trajectory: keep Shibboleth, add Keycloak local
%   D19.4 exit criteria reference Identity conformance tests (D16.5)
% ============================================================================

\chapter{Wave~1 --- Identity and federation}
\label{ch:wave1-identity}

\begin{caveatbox}[Dependencies]
\noindent Wave~1 applies the wave chapter template established in
\trchap{ch:wave0-foundations}{}. Its pre-requisite is Wave~0
(governance and inventory in place). It instantiates the identity
contract of \trchap{ch:identity}{}, and its closure is gated by
the four Identity tests of the conformance suite of
\trchap{ch:conformance}{} (\S\ref{sec:cf-suite}).
\end{caveatbox}

Wave~1 is, in most institutions, the wave that produces the most
visible operational change. The identity layer is touched by
nearly every institutional service; a wave that succeeds here
demonstrates that the playbook is operationally feasible, and a
wave that struggles here typically reveals limitations in the
governance established by Wave~0. The wave is consequently both
critical and instructive.

%-----------------------------------------------------------------------------
\section{Goal and dependencies}
\label{sec:w1-goal}

The goal of Wave~1 is to bring the institution's identity layer
into conformance with the contract of \trchap{ch:identity}{}: a
documented identity provider speaking standard interfaces (OIDC,
SAML2, SCIM~2.0), a federation participation aligned with the
applicable regional research-and-education identity federation,
an attribute schema reviewable by
the institution, and an authentication and assertion log
ingestible by the observability layer.

\paragraph{What the wave does not undertake.} Wave~1 does not
introduce policy-as-code (Wave~4), does not refactor endpoint
posture verification (Wave~6), and does not extend authentication
to workloads beyond what the identity layer naturally supports.
These are deferred to their respective waves; Wave~1's success
criterion is the contract, not the wider activation logic.

%-----------------------------------------------------------------------------
\section{Pre-requisites}
\label{sec:w1-prereq}

Wave~1 has Wave~0 as its only required pre-requisite, but four
operational conditions strengthen the wave when present:

\begin{itemize}[leftmargin=1.5em,nosep]
  \item the existing identity provider is inventoried, with its
    federation contracts, attribute schema, and downstream
    application set documented;
  \item the institution's source-of-truth for staff and student
    populations is identified, with the lifecycle events
    (creation, role change, departure) understood operationally;
  \item federation-policy responsibilities are clarified --- the
    institution knows which contracts it has signed with which
    federations and what attribute releases those contracts
    obligate;
  \item the network attachment of \trchap{ch:network}{} is in a
    stable state (not being actively migrated under Wave~3),
    since identity changes that interact with network attachment
    are best made against a stable substrate.
\end{itemize}

%-----------------------------------------------------------------------------
\section{Coexistence pattern}
\label{sec:w1-coexistence}

The principal engineering effort of Wave~1 is the coexistence of
the legacy and the incoming identity providers during the migration
period. Indicative durations are twelve to twenty-four months for
brownfield trajectories and substantially shorter for greenfield;
the actual duration is determined by the application portfolio's
re-federation pace.

\paragraph{Federation bridge.} The two identity providers are
brought into a federation relationship: each issues assertions
that the other can verify, applications can choose which provider
to trust during the migration, and a session begun on one is
recognisable to the other. The bridge is built on standard
interfaces (SAML2 metadata exchange, OIDC discovery), so that
neither provider becomes captive to the bridge configuration.

\paragraph{Application migration cohorts.} Applications are
migrated to the new provider in cohorts, typically organised by
risk profile (least-critical first), by organisational unit (one
unit at a
time for distributed institutions), or by technical similarity
(all OIDC-native applications, then SAML2-only). Each cohort's
migration is reversible within an agreed window; the next cohort
begins only when the previous cohort's reversibility window has
passed without incident.

%-----------------------------------------------------------------------------
\section{Trajectory: greenfield}
\label{sec:w1-greenfield}

For institutions in a greenfield position --- no significant
incumbent identity provider, or a deliberate decision to refound
the identity layer ---  Wave~1 is the most direct of the three
trajectories. The reference instantiation of
\S\ref{sec:id-reference} (an open-source identity provider with
SCIM provisioning, joined to the applicable regional
research-and-education identity federation) is built
without the burden of coexistence; federation participation is
established once; the application portfolio is integrated as it
is built. Indicative duration is on the order of six to nine
months.

%-----------------------------------------------------------------------------
\section{Trajectory: brownfield on an incumbent proprietary
identity suite}
\label{sec:w1-brownfield-ms}

For a common brownfield position --- an institution whose
identity layer already runs on \loc{incumbent-identity-suite-vendor}
--- Wave~1 is the wave
in which the institution's relationship with that incumbent
proprietary identity suite
is restructured rather than replaced. The architectural posture of
the playbook is to retain the incumbent suite as the identity
provider for
operational reasons (skills market, integration with the
institution's existing office-productivity estate) while adding the
infrastructure that prevents that suite from becoming the
institution's
policy plane in addition to its identity plane.

\paragraph{Pattern: incumbent suite retained, external policy
engine added in parallel.} The
institution retains the incumbent proprietary identity suite for
authentication, federation
participation, and SCIM provisioning. A new policy plane (Wave~4)
is prepared in parallel, which will eventually take over the
authorisation logic that the suite's proprietary conditional-access
expressions currently
encode. During Wave~1, that conditional-access configuration is
documented, simplified where possible, and pinned to expressions
whose translation to an open declarative policy language is known
to be
tractable. The objective is to bring Wave~1's exit criteria within
reach \emph{without} blocking on Wave~4.

\paragraph{What changes in Wave~1.} The federation metadata is
audited and, where necessary, brought into the applicable regional
research-and-education identity federation
in conformance with that federation's technical profile. The
attribute schema released to downstream applications is reviewed
and, where appropriate, reduced. The authentication and assertion
logs are routed to the observability layer through OpenTelemetry
collectors deployed in this wave. SCIM provisioning is verified
end-to-end against an alternative IdP test environment, to
demonstrate that exit-data export works.

\paragraph{What does not change.} The conditional-access logic
itself remains in place; it will be partially reproduced in the
external policy engine
during Wave~4. The office-productivity estate is unaffected. The
institution's existing procurement framework for the incumbent
suite continues
through this wave.

\paragraph{Indicative duration.} Twelve to eighteen months,
dominated by the application portfolio's federation review and the
SCIM provisioning verification.

%-----------------------------------------------------------------------------
\section{Trajectory: brownfield hybrid}
\label{sec:w1-brownfield-hybrid}

For research-intensive institutions with a mature, SAML-based
inter-institutional federation stack (for example a Shibboleth
deployment), Wave~1 is structured around the preservation of
that investment. The architectural posture is that the existing
federation stack
remains authoritative for inter-institutional federation, and a
local OIDC identity broker (for example a Keycloak instance) is
introduced for local role activation, attribute
enrichment, and the TAC session model that subsequent waves will
exercise.

\paragraph{Pattern: federation stack retained, local OIDC broker
added for local
activation.} Federation participation continues through the
existing federation deployment, with its
applicable-federation metadata,
attribute-release commitments, and operational team unchanged. A
new local broker instance is deployed as a component that
consumes the federation's assertions, enriches them with
institution-specific attributes, and emits short-lived OIDC tokens
that downstream applications consume.

\paragraph{What changes in Wave~1.} The local OIDC broker is
provisioned and
integrated with the existing source-of-truth via SCIM; the
authentication flow ``federation stack $\rightarrow$ local broker
$\rightarrow$
application'' is established for a first cohort of applications;
the observability layer receives events from both the federation
stack and the local broker. The federation contracts are not
renegotiated; the
metadata is not reissued.

\paragraph{What does not change.} Inter-institutional federation
behaviour, registration with the applicable identity federation,
and the relationships with peer
institutions remain unaffected. The applications that have not
yet adopted the local broker continue to authenticate directly
against the existing federation stack.

\paragraph{Indicative duration.} Nine to fifteen months. Shorter
than the incumbent-proprietary-suite trajectory because the
federation work
is largely already done; longer than the greenfield trajectory
because the dual stack must be operated with discipline.

%-----------------------------------------------------------------------------
\section{Reversibility criteria}
\label{sec:w1-reversibility}

Wave~1 is reversible at acceptable cost throughout the application
migration period and for an additional grace period after the last
cohort migrates.

\begin{itemize}[leftmargin=1.5em,nosep]
  \item The legacy identity provider is kept warm --- operational
    but with no new applications onboarding to it --- for an
    indicative period of six to twelve months after the final
    cohort migration.
  \item Per-cohort rollback is the unit of reversal: a cohort that
    encounters difficulties is rolled back to the legacy provider
    while the diagnostic proceeds; only when the diagnostic is
    closed is the next cohort moved.
  \item The full-wave rollback procedure is documented at wave
    open and exercised in a tabletop manner before any production
    cohort migration.
\end{itemize}

%-----------------------------------------------------------------------------
\section{Exit criteria}
\label{sec:w1-exit}

Wave~1 closes when the four Identity tests \emph{and} the three
Crypto-layer tests of the conformance suite of
\trchap{ch:conformance}{} (\S\ref{sec:cf-suite}) pass on the new
identity instantiation --- certificate and key services are stood up
here as the cryptographic prerequisite that identity depends on ---
when at least the agreed
percentage of applications has been migrated, and when the
governance body has reviewed the wave outcome.

\begin{itemize}[leftmargin=1.5em,nosep]
  \item OIDC discovery endpoint conformant; attribute schema
    institutionally controlled; SCIM export-import on alternative
    IdP preserves attributes; authentication and session events
    forensically reconstructable. (Identity conformance subset.)
  \item ACME/RFC~8555 certificate issuance conforms against the
    institutional CA; key-escrow recovery executes on a representative
    escrowed key; HSM key migration preserves cryptographic identity.
    (Crypto conformance subset --- the cryptographic foundation that
    identity is built on.)
  \item An institutionally agreed proportion of applications has been
    migrated, with the remainder either
    in a documented schedule or in the documented exception
    register.
  \item The governance body has reviewed the wave outcome and
    formally closed it, releasing the resources for the next
    waves.
\end{itemize}

\begin{realworldbox}[Identity migration precedents]
\noindent Several research-intensive institutions are reported to
have undertaken multi-year identity migrations of comparable scope.
Among the publicly discussed: CERN's migration to an open-source
SSO platform, ETH~Z\"urich's federation evolution, and various
campus
deployments described in the research-and-education networking
community (including national research and education network and
EDUCAUSE community fora).
Specific durations, application coverage, and operational outcomes
should be confirmed against current public documentation before
being cited; the patterns reported above are consistent with the
trajectories described in this chapter.
\end{realworldbox}

%-----------------------------------------------------------------------------
\section{Regulatory anchoring}
\label{sec:w1-regulatory}

\noindent Wave~1 establishes the identity-and-federation
substrate. It contributes to \gls{gdpr}~Art.~32(1)(b)
(access-control effectiveness) and Art.~25 (data minimisation
through controlled attribute release), and operationalises
NIS2~Art.~21(2)(i) (HR security and access control) and (j)
(MFA and secured communications). It also operationalises
eIDAS~2 trust services and notified-eID interoperability for
institutions that participate in the EUDI Wallet ecosystem. The
wave is the principal anchor of 27001~A.5.16, A.8.2 and A.8.5
and contributes to NIST CSF \textsf{PROTECT} (PR.AA --- identity
management, authentication and access control). Detailed mapping
in \trchap{ch:regulatory-mapping}{} \S\ref{sec:rm-gdpr},
\S\ref{sec:rm-nis2}, \S\ref{sec:rm-eidas2},
\S\ref{sec:rm-iso}, \S\ref{sec:rm-nist-csf}.
